\documentclass[answers,12pt,addpoints]{exam} 
\usepackage{import}

\import{C:/Users/prana/OneDrive/Desktop/MathNotes}{style.tex}

% Header
\newcommand{\name}{Pranav Tikkawar}
\newcommand{\course}{01:XXX:XXX}
\newcommand{\assignment}{Homework n}
\author{\name}
\title{\course \ - \assignment}

\begin{document}

I have done more literature review and have boiled the model down to three candidates. 

Our requirements are as follows:
\begin{itemize}
    \item Spatiotemporal model (in 1D for now)
    \item Isotropic in Space
    \item Non-Separable in any dimension
\end{itemize}
Ideal Conditions are:
\begin{itemize}
    \item Matern in space
    \item Matern/diffuse in time
\end{itemize}

The base spatial model is as follows $(\kappa^2 - \Delta)^{\alpha/2} Z(h) = W(h)$ where $\Delta$ is the laplacian, $\kappa$ is the range parameter, $\alpha$ is the smoothness paramter, $h$ is our coordinate (space), $Z$ is our process, and $W$ is a white noise process 

The first candidate for a spatiotemporal model is extending the laplacian to time 
$$ \Delta = \frac{\partial^2}{\partial t^2} + \sum_{x_i \in x} \frac{\partial^2}{\partial x_i^2} $$
This does not fit our critera but is good

Next is one where we add a first partial of time in the linear operator, $L = \kappa^2 + \Delta_s + \partial_t$ and then take $L^{\alpha/2}Z(h,u) = W(h,u)$ 
While I think this is a good start, I am doubtful on the flexibilithy of the model even when a sufficient amount of parameters are added. 

Another Candiadate is the model defined in the Paper https://www.stat.purdue.edu/~tlzhang/maternasym.pdf Where we consider the NFSST Matern Model 
$$ M_{r, a_1, a_2, \nu_1, \nu_2}(h, u) = E\left[\exp\left( -\frac{1}{2} \left( \frac{a_1^2 ||h||^2 }{V_1} +  \frac{2 a_1 a_2 u r^T h }{ \sqrt{V_1 V_2}} + \frac{a_2^2 u^2 }{V_2}   \right) \right) \right]$$
Where the $h$ is the spacial coordinate, $u$ is the temporal coordinate.

$M$ is called the NFSST (Not fully symmetric space time) Matern Correlation function. 

I really like this model. Since it fully fits every single requirement and ideal condition. Also I can see a great next step being a SPDE based FEM solver since I am not able to find much done on this specific topic. 

Next steps are as follows 
\begin{itemize}
    \item Review the Paper further, especially what the interaction parameter is and how to estimate it. 
    \item Understand the fourier transform of this model for the frequency behavior and then the PIT of the model both in space and time.
    \item Develop a linear operator $L$ s.t. it has the same spectral density as the model. 
    \item Develop a method to solve the SPDE numerically using some FEM method and potentially adapting the R-INLA or RSPDE libraries solver.
\end{itemize}


\end{document}